\section{Spinor transmission and the residual chiral energy}

In an oriented Riemannian spin four-manifold let $P$ be the twistor
operator on positive half-spinors. We use the curvature operator with
positive sectional curvature on the round sphere, the Hilbert--Schmidt
norm on $W^+$, and the Hilbert norm on spinor-valued two-forms.

\begin{lemma}[Weighted finite cover]
\label{lem:finite-cover-defect-bound}
\lean{FirstOrderTwistor.finiteCoverDefectBound, FirstOrderTwistor.finiteCoverDefectBoundOfIntegrals}
\leanok
Let $r\le u+\sum_{j\in J} w_j$ for a finite $J$. If $c_j>0$ and
$c_j^2 w_j\le 2H_j$ for $j\in J$, then
$r\le u+\sum_{j\in J}2H_j/c_j^2$.
For a nonnegative, integrable density on a measurable residual $S$,
take $w_j$ to be its integral on $S\cap U_j$ and $u$ its integral on
$S\setminus\bigcup_{j\in J}U_j$, for measurable strips $U_j$.
\end{lemma}
\begin{proof}
\leanok
The energy measure restricted to $S$ is finite. Subadditivity on the
union gives the cover inequality even with overlaps. Divide each strip
estimate by its positive squared floor and sum. An empty family leaves
the whole residual in the uncovered term.
\end{proof}

\begin{proposition}[Projected twistor curl in a supplied normal frame]
\label{prop:supplied-normal-spin-defect}
\lean{FirstOrderTwistor.selfDualCliffordProjection_orthogonal, FirstOrderTwistor.spinCovariantDerivative_commutator, FirstOrderTwistor.spinConnectionLift_commutator, FirstOrderTwistor.selfDualCurvatureBlock_metric_isSymm, FirstOrderTwistor.metricProjectedFrameTwistorCurl_normSq, FirstOrderTwistor.metricCoordinateTwistorDefect, FirstOrderTwistor.spinNormalGauge_leviCivita_at_center, FirstOrderTwistor.metricNormalChartTwistorDefect}
\leanok
In the positive quaternionic half-spin fiber the unit self-dual forms
act by $C_a=\sqrt{2}\,I_a$ for $I_a=i,j,k$. Write $m$ for contraction
and $\Pi=1-m^*m/6$. In a coordinate model supply a $C^2$ field, a
differentiable frame connection vanishing at the point, upper and lower
Clifford symbols with their normal-frame values and zero first derivatives,
and an identification of the frame curvature with the metric curvature
in the chosen oriented orthonormal frame. For the coordinate twistor
operator and its exterior covariant derivative,
\[
 \Pi(d^\nabla P\psi)^+_a=-\tfrac12 c(W^+\omega_a)\psi,
 \qquad
 \sum_a|\Pi(d^\nabla P\psi)^+_a|^2
   =\tfrac12|W^+|_{\rm HS}^2|\psi|^2 .
\]
In an aligned boundaryless normal chart the hypotheses on the symbols and
curvature follow from the metric Christoffel matrices and a supplied
twice differentiable moving frame with identity value and zero first
derivative at the center and a differentiable local inverse. For spinor
interpretation choose an oriented orthonormal frame field and its positive
spin lift. The covector connection term vanishes at the normal center.
The squared formula also holds for every ordered two-jet realizing
the metric spin curvature commutator in that frame.
\end{proposition}
\begin{proof}
\leanok
The commutator of two coordinate covariant derivatives equals the
curvature of the supplied connection by symmetry of the ordinary
Hessian. The positive spin lift preserves brackets of skew matrices.
Self-dual alternation of the Clifford correction in $P$ lies in the
image of $m^*$, as does the scalar part of the curvature block.
The trace-free symmetric block lies in $\ker m$, and $mm^*=6$.
The norm of left Clifford multiplication by a real self-dual form
squares to twice the squared form norm times the squared spinor norm.
\end{proof}

\begin{proposition}[Measured strip costs on the physical remainder]
\label{prop:spin-remainder-costs}
\uses{lem:finite-cover-defect-bound,prop:supplied-normal-spin-defect,thm:covered-measure-accounting}
\lean{FirstOrderTwistor.spinDefectFloorBound, FirstOrderTwistor.spinDefectEnergyDecomposition, FirstOrderTwistor.SpinDefectStrip.validOn_ofMetricJets, FirstOrderTwistor.RemainderEstimates.ofSpinDefectStrips, FirstOrderTwistor.RemainderEstimates.ofNormalizedSpinDefectStrips, FirstOrderTwistor.SpinDefectStrip.normalizedCost_le_of_singleton_lpNorm, FirstOrderTwistor.energyIdentityOfSpinorRatioBounds}
\leanok
On each strip intersected with the physical remainder, assume equality
of the restricted physical and strip chiral energy measures. In the
strip metric write $S_j=\sum_a|\psi_{ja}|^2$ and
$H_j=\sum_a|\Pi(d^\nabla P\psi_{ja})^+|^2$. Given the a.e. per-field
identities and the required integrability, a joint floor $S_j\ge c_j^2>0$
gives cost $2\int H_j/c_j^2$. Alternatively use the joint nonzero set
and cost $2\int H_j/S_j$ there. In the latter case the uncovered set
consists of points outside every contributing joint nonzero set. Each
strip uses its own fields and volume. If the summed costs and the uncovered
energy both vanish, first in the eventual sequence bounds at fixed
truncation and then on exhaustion, these estimates give the remainder
hypothesis of Theorem~\ref{thm:covered-measure-accounting}.
\end{proposition}
\begin{proof}
\leanok
Sum the identities over fields and integrate. On the positive locus the
ratio is exactly half the curvature density. On a floor set use its
lower bound in the unnormalized identity. Transfer only the restricted
energy measure. At each pair of indices take the finite union and retain
its complement once. The number of strips may vary with both indices.
\end{proof}

\begin{lemma}[First-order perturbations and a differentiated interior bound]
\label{lem:spin-strip-derivative}
\uses{lem:buffered-euclidean-interpolation}
\lean{FirstOrderTwistor.spinStrip_operatorPerturbation, FirstOrderTwistor.spinStrip_bufferedCoefficientBound, FirstOrderTwistor.spinStrip_sumSlabBounds, FirstOrderTwistor.spinStrip_absorbTwistedDirac, FirstOrderTwistor.spinStrip_projectedDerivative, FirstOrderTwistor.harmonicMetricCoordinateTwistedDiracDerivative_normSq}
\leanok
For a supplied operator difference $a\nabla+b$ on a reference strip,
gradient norm bound one and Sobolev constant $S$, the operator norm
from the first-order space to $L^2$ is at most
$\|a\|_\infty+S\|b\|_4$. For the interior estimate on $K\Subset K'$,
supply the flat cutoff estimate, small twisted-operator perturbation,
the differentiated harmonic equation with its cutoff and Ricci forcing
bounds, the return to the actual norms, an incoming error $d$, a
first-order comparison error $\epsilon_1$, and a floor $|\psi|\ge c>0$
on $K$. With constants from those estimates,
\[
 \left\|\frac{|\Pi(d^\nabla P\psi)^+|}{|\psi|}\right\|_{2,K}
 \le \frac{C}{c}\left(\epsilon_1+d+
       \|\,|\mathrm{Ric}|\,|\psi|\,\|_{2,K'}\right).
\]
\end{lemma}
\begin{proof}
\leanok
Apply H\"older with exponents $(\infty,2)$ and $(4,4)$ in the supplied
measure. Buffered interpolation bounds the coefficient budget under its
chart comparison hypotheses. For the derivative estimate, absorb the
relative error in the flat estimate when $C_0\epsilon\le1/2$. Bound the
cutoff term on the outer buffer, then divide the projected norm by the
floor on the inner set.
\end{proof}

\begin{proposition}[Relative Bochner criterion without a positive minimum]
\label{prop:spin-relative-bochner}
\lean{FirstOrderTwistor.normalizedSpinCutoff_fderiv_bound_frame, FirstOrderTwistor.spinRelativeBochner_bound, FirstOrderTwistor.spinMixedRelative_sum_le, FirstOrderTwistor.spinCutoffIntegral_in_gauge, FirstOrderTwistor.RemainderEstimates.ofSpinCutoffCover, FirstOrderTwistor.energyIdentityOfSpinBochnerBudgets}
\leanok
For measured joint spinor mass $S$, derivative density $Q$, and cutoff
$\chi$, set $v^2=Q/S$ on the joint nonzero locus. Assume $S>0$ almost
everywhere on the cutoff support. Given the integrated regularized
twisted Bochner inequality and finite relative costs in the strip metric,
\[
 \int\chi^2|W^+|^2 \le C\int
 \big(\chi^2|\mathrm{Ric}|^2+\chi^2|\mathrm{Rm}|v^2+
          \chi^2v^4+|d\chi|^2v^2\big).
\]
For a finite family of measured patches the sum of the mixed terms
is at most $\sqrt{\mathcal K\mathcal V}$, where $\mathcal K$ and
$\mathcal V$ are the sums of $\int\chi^2|\mathrm{Rm}|^2$ and
$\int\chi^2 v^4$ in the respective gauges. Where the cutoff squares
sum to at least one in the physical energy measure, the weighted costs
bound the covered energy. The mass of the remaining set enters once.
For a uniform constant, bounded $\mathcal K$, vanishing Ricci-square,
relative fourth-power and cutoff costs, and vanishing uncovered energy
in the iterated order, Route~02 gives the component energy identity
under its truncation and total convergence hypotheses.
\end{proposition}
\begin{proof}
\leanok
Writing $D=S+\epsilon$, differentiating $\chi/\sqrt D$ gives
$|d(\chi/\sqrt D)|^2Q\le 2|d\chi|^2 Q/D+2\chi^2(Q/D)^2$
in the Hilbert frame model. Use the factor $S/(S+\epsilon)$ on the left and dominate the regularized
right side by the finite unregularized costs. Dominated convergence
takes $\epsilon$ to zero. Apply Cauchy--Schwarz first to each integral,
then to the finite sum over patches.
\end{proof}

\begin{proposition}[Incoming modes and attenuation]
\label{prop:spin-mode-transmission}
\lean{FirstOrderTwistor.forcedCylinderMode_equation, FirstOrderTwistor.cylinderMode_energy_balance, FirstOrderTwistor.fixedDomainDiracPerturbation_inverse_norm_le, FirstOrderTwistor.fixedDomainDiracScattering_sameEnd_norm_le, FirstOrderTwistor.diracWeighted_remoteTrace_norm_le, FirstOrderTwistor.roundCylinderMidpoint_jointFloor_Lp}
\leanok
For $\partial_t+A$ on $[0,L]$, a positive signed mode uses its left
datum, a negative mode its right datum. On a supplied complete
orthonormal modal model with $|\lambda|\ge3/2$, a finite family with
joint floor $c^2$ on the middle slice satisfies
\[
 c^2\operatorname{vol}(Y)
 \le e^{-3L/2}\sum_a(\|x_a\|_2^2+\|y_a\|_2^2)
\]
for the homogeneous solutions and their signed incoming components.
For a boundedly invertible reference operator with inverse norm at most
$C$, a fixed-domain perturbation of norm at most $\delta$ has inverse
norm at most $C/(1-C\delta)$ if $C\delta<1$.
With fixed incoming and outgoing trace maps and a vanishing same-end
reference block, the resolvent identity bounds the same-end reflection
by the perturbation norm times the two inverse and trace factors.
\end{proposition}
\begin{proof}
\leanok
Integrate each signed scalar equation from its incoming end. At distance
$L/2$ the squared amplitude decreases by at least the indicated exponential
factor. Sum using the Hilbert norm and integrate the floor on the slice.
The perturbation bound follows by a Neumann inverse on the fixed spaces.
\end{proof}

\begin{proposition}[Matching and regular jump corrections]
\label{prop:spin-core-and-jump}
\lean{FirstOrderTwistor.matchSuppliedCoreTraceGraphs, FirstOrderTwistor.closedJumpCorrectionOfRightInverse, FirstOrderTwistor.criticalJointFloorOfJumpBounds}
\leanok
For two supplied bounded core graphs $G_L,G_R$ and propagation norm
bounds $e^{-3L/2}$ in both directions, the affine graph equations have
a unique solution when $\|G_L\|\|G_R\|e^{-3L}<1$. For an original-source
jump problem supply four measured face maps, a right lift of both jumps,
a compatible closed inverse, and zero-jump gluing. There is a unique
piecewise homogeneous correction with the specified jumps. If its regular
fields satisfy $\|\psi_a\|_{8/3}\le C J_a$ and their joint floor on $U$ is
$c^2$, then $c^2\operatorname{vol}(U)^{3/4}\le C^2\sum_a J_a^2$ for
nonnegative jumps and bound constants as in the critical norm estimate.
\end{proposition}
\begin{proof}
\leanok
Substitute one affine trace equation into the other and apply contraction.
For the jump correction subtract from the piecewise lift the restriction
of the global inverse applied to its Dirac error. Use the $L^{4/3}$
triangle inequality on the sum of squared field norms.
\end{proof}

\begin{proposition}[Attainment in a supplied compact gauge]
\label{prop:spin-compact-kernel}
\lean{FirstOrderTwistor.existsUnitKernelOfCompactEmbedding, FirstOrderTwistor.weylBlock_eq_zero_of_dense_spinor}
\leanok
Fix a real Hilbert domain with compact inclusion into the reference
$L^2$ space. A bounded sequence whose image norms tend to one, whose
operator defects tend to zero, and whose operators converge in norm to
$P$ yields a kernel element with unit image norm. Separately, a continuous
Weyl block vanishes throughout a chart if its projected defect vanishes
and the spinor is nonzero on a dense set satisfying the pointwise identity.
\end{proposition}
\begin{proof}
\leanok
Take a strongly converging image subsequence and a weak cluster point in
the bounded Hilbert domain. Continuous linear maps identify the kernel
and the image norm of this point. For the Weyl block, apply the pointwise
identity on the dense set and continuity to all coefficients.
\end{proof}

\begin{openproblem}[Differentiated control over an entire physical residual]
\label{problem:physical-spin-residual}
\notready
Construct measured fields or local gauges on the actual residual of the
geometric sequence with summably small independent differentiated bounds
and negligible uncovered energy. A Green pole, compatible endpoint germs,
a fixed-domain Sobolev inverse, and a normalized limit on the same
smooth manifold each require their own geometric realization.
\end{openproblem}
